\section{Future Directions}
\label{sec:future-directions}

The preceding sections show that open-vocabulary and reasoning segmentation have moved beyond closed-set dense prediction. CLIP-based models make category names and descriptions usable at test time; SAM-style models make high-quality masks accessible through prompts; grounded and MLLM-based systems further connect language, reasoning, and pixels~\citep{radford2021learning,kirillov2023sam,ravi2024sam2,xlaiLisaReasoningSegmentation2024,renPixellmPixelReasoning2024}. Nevertheless, the comparison in Section~\ref{sec:comparison-discussion} also shows that stronger interfaces create new weaknesses: benchmark protocols are fragmented, reasoning traces are not always faithful to masks, hybrid pipelines are difficult to reproduce, and domain extensions require more than ordinary mIoU. The following directions therefore follow from the limitations already discussed rather than from a separate list of speculative topics.

\subsection{Unified and transparent evaluation protocols}

Existing studies often report strong results under different vocabularies, prompt templates, mask generators, training data, foundation backbones, and metrics. This makes direct comparison difficult even when methods are evaluated on datasets with similar names, because a gain may come from improved category-name renovation, stronger proposal generation, additional reference data, or a different CLIP/SAM variant rather than from the core algorithm~\citep{liang2023ovseg,xu2023san,hhuangRenovatingNamesOpenvocabulary2024,ybenigmimFlossFreeLunch2025,shiyHarnessingVisionFoundation2025}. Future benchmarks should therefore require a transparent reporting checklist that includes class-name sources, synonym and prompt rules, open-vocabulary training exposure, prompt source, mask generator, foundation modules, compute, latency, and all post-processing steps.

For reasoning and grounded segmentation, unified evaluation also needs to record the input interface. A category list, a referring phrase, a grounded caption, a false-premise query, a medical reasoning instruction, and a video question do not test the same capability~\citep{chngMaskGroundingReferring2024,rasheedGlammPixelGrounding2024,huGroundingSuiteMeasuringComplex2025,yanMedreasonerReinforcementLearning2026,zhengVillaVideoReasoning2025}. More research is expected on benchmark suites that report mask quality together with target-selection accuracy, answer correctness, rejection accuracy, and failure categories, so that a high mask score cannot hide an incorrect language decision.

\subsection{Faithful reasoning and verifiable pixel grounding}

Reasoning segmentation is valuable because it can infer implicit targets, explain object selection, and output masks in one interaction. However, current systems can still produce plausible explanations that are not causally tied to the mask, or accurate masks whose textual rationales rely on the wrong evidence~\citep{xlaiLisaReasoningSegmentation2024,renPixellmPixelReasoning2024,zhuLensLearningSegment2026,zhouReasoningImplicitSelfsupervised2026}. This issue is important because the intended use of reasoning segmentation is not only to draw a mask, but also to make the mask accountable to a user query.

Future work should make the relation among query, reasoning trace, target grounding, and final mask explicitly checkable. False-premise handling, rejection tokens, backward grounding verification, and grounded visual chat already point in this direction~\citep{xiaGsvaGeneralizedSegmentation2024,wuSeeSaySegment2024,zhangLlavagroundingGroundedVisual2024,guoSeeingBelievingRichContext2026}. More efforts are needed on causal grounding tests, counterfactual queries, explanation-mask consistency metrics, and training objectives that reward correct target inference rather than only final IoU.

\subsection{Efficient and reproducible hybrid segmentation systems}

Many recent methods obtain strong performance by composing several foundation models. Trident combines CLIP, DINO, and SAM; CLIPer uses CLIP spatial repair and diffusion-based compensation; Grounded-SAM pipelines attach detectors, language grounding models, agents, or MLLMs to SAM/SAM2~\citep{shiyHarnessingVisionFoundation2025,lsunCliperHierarchicallyImproving2025,liu2024groundingdino,luoConnectingDotsTrainingFree2026}. These systems are attractive because they avoid full retraining, but they also make cost, latency, memory, and reproducibility harder to interpret.

Future research should study hybrid systems as deployable systems rather than isolated algorithms. Lightweight adapters, segment embeddings, mask-aware CLIP adaptation, compact pixel decoders, and reward-optimized segmentation tokens provide possible routes for reducing the dependence on large tool chains~\citep{xu2023san,yliMaskadapterDevilMasks2025,xwangUseUniversalSegment2024,renPixellmPixelReasoning2024,huangSamr1LeveragingSam2026}. Given the limited comparability of current reports, future papers should publish module versions, prompt budgets, inference calls, runtime, memory, and failure cases together with accuracy.

\subsection{Domain-robust segmentation beyond ordinary 2D images}

The extension from natural images to video, 3D, remote sensing, medical images, and embodied or affordance settings changes the meaning of segmentation. Video methods must preserve temporal identity and motion consistency; 3D methods must ground language in geometry and viewpoints; remote-sensing methods must handle scale, rotation, and sensor variation; medical methods must follow clinical target definitions; affordance methods must segment functional regions rather than whole objects~\citep{baiOneTokenSeg2024,zhengVillaVideoReasoning2025,chen3DDRESDetailed3D2026,huangSurprise3dDatasetSpatial2026,libExploringEfficientOpenvocabulary2026,yanMedreasonerReinforcementLearning2026,wangAffordancer1ReinforcementLearning2026}. Treating all of these settings as ordinary image segmentation would erase the very constraints that make them difficult.

Future studies should design domain-specific protocols while preserving a shared reporting backbone. For example, video reasoning segmentation should jointly report answer accuracy and $\mathcal{J}\&\mathcal{F}$; clinical grounding should report target-interpretation errors separately from mask Dice; affordance grounding should distinguish object localization from functional-part localization~\citep{zhengVillaVideoReasoning2025,yanMedreasonerReinforcementLearning2026,wangAffordancer1ReinforcementLearning2026}. More research is expected on benchmarks and models that keep these domain constraints explicit, rather than hiding them behind a single aggregate score.

\subsection{Human-in-the-loop and accountable segmentation}

Open-vocabulary and reasoning segmentation are increasingly used through interactive prompts, conversations, and correction loops. Yet most benchmarks still evaluate one-shot predictions, while real users often refine targets, correct masks, change instructions, or ask why a region was selected~\citep{kirillov2023sam,ravi2024sam2,rasheedGlammPixelGrounding2024,wuSeeSaySegment2024}. This mismatch matters because a system that is imperfect in one shot may still be useful if it can accept corrections reliably, and a system with high one-shot accuracy may still be unsafe if it cannot explain or reject uncertain requests.

Future work can therefore combine promptable segmentation, grounded dialogue, active correction, and uncertainty reporting. The key research problems include how to update masks without losing open-vocabulary semantics, how to preserve conversation-state consistency across turns and frames, how to expose uncertainty to users, and how to audit final masks in high-stakes domains. This direction naturally connects the technical progress reviewed in Sections~\ref{sec:promptable-grounded-segmentation}--\ref{sec:extended-scenarios} with the evaluation requirements summarized in Section~\ref{sec:comparison-discussion}.
